\chapter{预训练 VLM 从哪里来，它究竟在计算什么}

\section{先分清三个时间阶段}

你现在使用的 \texttt{Qwen3-VL-2B-Instruct} 是已经完成通用训练、能够接受图像和文字的模型。本地微调从这些已有权重出发。为了理解后续每个文件和命令，先把时间轴分成三段：

\begin{enumerate}
\item \textbf{厂商训练基础模型。} 大量通用数据决定它怎样把图像与语言联系起来，以及怎样理解问题、生成回答。这些训练发生在我们下载模型之前。
\item \textbf{我们做机器人任务微调。} 使用相机画面、语言任务、状态和示范动作，让已有模型适应小鸭子的技能定义及执行条件。这是当前正在进行的 QLoRA 监督微调。
\item \textbf{部署后进行推理。} 把新画面和新指令交给模型，生成一次决策，执行后再观察。普通推理不会更新权重；连续走了十步，也不意味着自动训练了十次。
\end{enumerate}

这里的“预训练模型”是从使用者角度说的：它是我们任务训练的起点。具体下载的 Instruct 检查点还包含厂商的指令后训练，不能把它理解为只学过下一词预测、完全没学过对话的原始模型。官方提供 Instruct 与 Thinking 等版本；本项目使用前者。\cite{qwen-card}

\begin{notebox}{一个必须明确的边界}
会识别“球”、理解“红色”和“向左”，不等于知道 HD1910 的每个关节该转多少，也不等于知道本项目 \texttt{kick\_left} 的参数和前置条件。技能库由本项目提供；底层动作由已有控制器执行。通用视觉语言知识、上层技能决策、底层运动控制是三个相互配合的部分。
\end{notebox}

\section{模型名称、文件与参数分别表示什么}

\texttt{Qwen3-VL} 是模型家族；VL 指视觉与语言；\texttt{2B} 是规模档位名称，不能简单地用“2B 乘每参数字节数”代替完整显存预算；\texttt{Instruct} 表示面向指令交互的版本。模型卡将它归类为图像加文本到文本的模型，并标注 Apache-2.0 许可证。本地调用无需按每次请求向 API 服务商付费，但设备、存储和电力仍有成本；采用云端服务时计费由具体服务决定。\cite{qwen-card}

“权重”是一大组训练得到的数值，保存着模型所学的统计规律。“配置”规定层数、维度等结构；“Tokenizer”规定文字如何分成 token；“Processor”还处理图片及聊天模板。因此，仅有一个权重文件并不等于已经拥有正确的完整输入管线。

本地下载的配置与官方 2B 配置一致：语言部分有 28 层、隐藏维度 2,048、词表大小 151,936；视觉部分有 24 层、隐藏维度 1,024，输入为 3 个颜色通道，图像 patch 大小为 16，空间合并系数为 2。配置中的最大位置长度为 262,144；这是结构支持值，不代表本机能按这个长度训练。本次本地训练的上下文上限是 4,096。\cite{qwen-config}

\section{从一张图片到一句回答}

\subsection{第一步：把人类输入整理为消息}

人类看到的是“这张图里有什么？”和一张图片。程序看到的是一组有角色的消息：\texttt{system} 给出任务约定，\texttt{user} 提供图像和问题，\texttt{assistant} 是模型要生成的回答。一个用户消息可以包含多个图像块和文本块。Processor 负责聊天模板、文字分词和图像预处理。\cite{multimodal-template}

下面是教学用伪结构，帮助理解数据方向，并不是复制即可运行的模型调用代码：

\begin{Verbatim}[breaklines=true,fontsize=\small]
messages = [
  {role: "system", content: [task_and_output_rules]},
  {role: "user", content: [front_rgb, near_rgb, instruction, state]}
]
processed = processor(messages)
answer_tokens = model.generate(processed)
answer_text = tokenizer.decode(answer_tokens)
\end{Verbatim}

训练时还会附加已知的正确 assistant 回答；推理时只给到 assistant 应该开始回答的位置。两者输入格式要对齐，否则模型训练时见到的分布与实际执行时不同。

\subsection{第二步：图片变成视觉特征}

模型并不是先把图片完整地翻译为一句中文再思考。它会把像素整理为小块，让视觉编码器计算一组数值向量。例如，一个向量可能携带局部颜色、边缘、纹理与物体关系等混合信息；不能把每个维度都直接解释成一个人类概念。

Qwen3-VL 的主要组成是视觉编码器、视觉与语言之间的合并／投影模块以及语言模型。系列报告说明，2B 和 4B 使用 SigLIP2-Large 作为视觉编码器基础；视觉特征经空间合并后进入语言模型。\cite{qwen-report} 合并与投影的作用，可以理解为把视觉侧的数值表示变成语言侧能利用的表示。维度匹配只是必要条件，跨模态训练还要让这些表示具有可用的对应关系。

“视觉 token”与“文字 token”不能混为一谈。文字 token 有词表编号，可以解码为文字片段；视觉 token 通常是从图片计算出的向量序列。二者在多模态上下文中相互作用，并不表示每块图片都必定对应某个汉字。

Qwen3-VL 还通过 DeepStack 将不同深度的视觉特征融入语言层，使用交错式 MRoPE 表达空间／时间位置，并通过文字时间戳处理视频时间关系。这些是模型家族的结构能力。\cite{qwen-repo} 当前小鸭子每次提交的是两张观察图片，不能因为基础模型支持视频，就把当前管线说成已经输入连续视频。

\subsection{第三步：数值张量进入模型}

典型预处理结果包含以下几类对象。它们不是四份独立答案，而是一次模型计算需要的输入。\cite{multimodal-template}

\begin{longtable}{P{4.0cm}P{10.5cm}}
\textbf{名称} & \textbf{通俗含义} \\
\hline
\texttt{input\_ids} & 文字及特殊占位标记的整数序列；每个整数是一个 token 编号。\\
\texttt{attention\_mask} & 告诉模型哪些位置有效，哪些只是批处理填充。\\
\texttt{pixel\_values} & 图片经缩放、归一化、分块等处理后的像素张量。\\
\texttt{image\_grid\_thw} & 视觉网格的时间、高、宽结构，供位置及视觉序列组织使用。\\
\end{longtable}

同一张图片的原始像素、视觉特征、最后输出的 token 是三个不同层级。原始 RGB 通常是三通道颜色值；隐藏状态是一组浮点向量；最终输出则是可以解码为文字的词表编号。图片分辨率、图片数量、提示词和历史长度都会改变一次计算的工作量。

\subsection{第四步：逐个生成输出 token}

语言模型根据当前全部输入以及已经生成的回答前缀，计算下一个 token 的分数，归一化后得到概率分布。选择一个 token、接到回答后面，再继续下一次预测，直到结束标记或达到生成长度限制。官方推理接口最终把生成的编号解码为文本。\cite{qwen-transformers}

如果回答是“前面有一个红球”，模型并不是从固定句子表里一次挑中整句；如果回答是 JSON，它同样要逐段生成括号、键名、值和标点。一个动作名称可能由多个 token 组成，不能假设 \texttt{move\_forward} 必定只占一个 token。

\section{预训练到底怎样让模型学会看图说话}

\subsection{用一个简化样本理解学习过程}

设训练样本包含“一张桌面上放着红球的图片”，以及目标描述“桌面上有一个红色的球”。这只是自造的教学示例，不是对官方真实训练样本的转述。

训练会让模型在看到图片和前文后，依次预测正确描述中的 token。假设应输出“红”，但模型给“蓝”的概率更高，损失就会较大。反向传播计算哪些可训练参数使这个错误更严重，优化器再小幅调整它们。很多不同图片、不同语言表达、不同问题反复出现后，模型逐渐建立可迁移的视觉与语言联系。

用公式表示一种典型自回归目标：
\[
  L=-\frac{1}{T}\sum_{t=1}^{T}\log p_{\theta}
  \left(y_t\mid x,y_1,\ldots,y_{t-1}\right).
\]
其中，$x$ 是图像和文字条件，$y_t$ 是第 $t$ 个目标 token，$\theta$ 是当前可训练参数，$T$ 是参与计算损失的目标长度。正确 token 的概率越高，对应损失越低。这是用于理解的基础表达式；官方大规模训练可能采用不同样本权重、数据配比和更复杂目标，不能把这个简式当成其完整复现配方。

监督训练通常使用正确的回答前缀，即 teacher forcing。训练预测到“红色”时，前面已经给出正确文字；自由生成时，前文是模型自己刚生成的文字。于是，“正确前缀条件下的验证损失降低”与“自己完整生成的回答可靠”并不等价。这也是我们既要计算验证损失，又要做自由生成和仿真闭环测试的原因。

\subsection{官方公开了哪些阶段}

下面仅压缩列出系列技术报告公开的预训练设置，B 表示十亿、T 表示万亿 token；token 预算不是图片张数。\cite{qwen-report}

\begin{longtable}{P{1.1cm}P{4.0cm}P{3.0cm}P{2.4cm}P{2.4cm}}
\textbf{阶段} & \textbf{目的} & \textbf{更新部分} & \textbf{token 预算} & \textbf{序列长度}\\
\hline
S0 & 视觉语言对齐 & 合并／投影层 & 67B & 8,192\\
S1 & 多模态预训练 & 全部组件 & 约 1T & 8,192\\
S2 & 长上下文学习 & 全部组件 & 约 1T & 32,768\\
S3 & 超长上下文适配 & 全部组件 & 100B & 262,144\\
\end{longtable}

对齐阶段可理解为先让已有视觉能力和语言能力连得起来；后面的联合训练进一步调整它们的协作方式；更长上下文训练让模型有机会学习跨度更大的依赖。这个解释说明每个阶段要解决的问题，并不意味着你的 6GB 显卡需要重做厂商的全部流程。

\subsection{用过什么数据，哪些细节仍不知道}

系列报告列举的类型包括图像描述、图文交错网页／文档、视觉问答、OCR、定位与计数、空间与三维标注、代码、视频、GUI／工具使用，以及纯文本。定位数据举例包含 COCO、Objects365、OpenImages、RefCOCO 和 PixMo；数据还包含内部采集及合成内容。\cite{qwen-report}

可以用下面这些\textbf{自造教学样本}理解各类监督信号的差别：

\begin{longtable}{P{3.0cm}P{6.0cm}P{5.3cm}}
\textbf{任务类型} & \textbf{输入示例} & \textbf{目标输出示例}\\
\hline
图像描述 & 一张桌面照片 & “桌上有红球和蓝杯子。”\\
视觉问答 & 图片＋“球是什么颜色？” & “红色。”\\
文字识别 & 一张标牌照片 & 标牌上可见的文字。\\
视觉定位 & 图片＋“指出红球的位置” & 目标点或边界框坐标。\\
关系判断 & 图片＋“球在杯子的哪一侧？” & “在杯子的左侧。”\\
视频理解 & 一段带时间信息的画面＋问题 & 事件描述或对应时间段。\\
工具调用 & 观察＋任务＋可用工具说明 & 某个工具名称及参数。\\
机器人技能监督 & 两路图＋任务＋状态＋历史 & 本项目下一步技能的 JSON。\\
\end{longtable}

前几类能力能帮助理解机器人场景，但最后一行仍需要本项目定义。即使通用数据中有工具调用任务，模型也不能凭空获知一套尚未展示给它的私有工具接口。

\begin{notebox}{公开资料的范围}
官方系列报告不是当前 2B-Instruct 检查点逐样本、逐轮次的训练清单。公开的类别和阶段设置不能反推出每个小模型见过的所有样本、精确比例、随机顺序或全部过滤阈值。因此不能声称它预训练时看过我们的 HD1910、已学过我们的踢球控制器，或给出没有公开依据的完整数据配方。
\end{notebox}

\section{Instruct 是怎样在预训练后形成的}

系列技术报告的后训练流程包含监督微调（SFT）、强模型到弱模型的知识蒸馏以及强化学习（RL）；标准回复与显式推理格式的数据面向不同版本。\cite{qwen-report} 这些概念分别解决不同问题：

\begin{itemize}
\item \textbf{SFT：学习示范。} 给定问题和目标回答，用目标回答监督模型。你当前的机器人训练也采用这一基本形式，只是示范来自任务规则和仿真记录。
\item \textbf{蒸馏：学习更强教师的行为。} 教师可以提供答案或概率信息，学生学习模仿其中有用的能力。教师是规则程序时，通常更直接称为示范模仿；不应自动宣称使用了厂商相同的蒸馏算法。
\item \textbf{RL：依据奖励改变策略。} 训练中生成候选行为，依据可验证结果或其他奖励信号优化。保存失败案例本身不是 RL；必须把奖励与参数更新过程实际连接起来。
\end{itemize}

厂商用过 RL，不意味着我们现在也在跑 RL；我们接着一个经历过多阶段训练的基础模型，用自己的监督样本进行 QLoRA。SFT 描述学习信号，QLoRA 描述参数更新与存储方式，两者可以同时成立。

\section{它的输出格式究竟由谁决定}

基础模型的自然输出是文本 token 序列，没有天生绑定我们五字段的机器人 JSON。固定格式来自几层约束的共同作用：

\begin{enumerate}
\item \textbf{系统提示} 告诉模型可用技能、各参数的意义，以及只能生成什么格式。
\item \textbf{微调样本} 反复示范相同格式和合理决策，使输出更符合任务。
\item \textbf{解码约束（如果实际启用）} 限制可生成的语法结构；本次不能把“提示要求 JSON”自动说成实现了强制语法解码。
\item \textbf{程序解析和守卫} 检查 JSON 字段、范围和当前执行条件。格式正确仍可能选错技能，因此执行前还要检查动作是否允许。
\end{enumerate}

例如，下列是本项目协议的一种教学示例，而不是该基础模型原生保证的输出：

\begin{Verbatim}[breaklines=true,fontsize=\small]
{
  "type": "skill",
  "skill": "move_forward",
  "seconds": 0.2,
  "target_color": "red",
  "reply": ""
}
\end{Verbatim}

“红色”来自用户目标与场景观察之间的关联：指令确定想找哪种球，图像及测量确定哪一个球在哪里。技能定义来自程序提供的提示与训练示范。输出格式来自项目协议，模型学习遵守它；最终是否执行由程序决定。

这一结构当前最清楚的叫法是“VLM 驱动的技能调用”。经过机器人动作示范训练后，也可按动作层级称为“分层视觉语言动作策略”。名称不是验收指标：真正要检查的是输入里有什么、输出是什么、哪些参数被更新，以及闭环任务是否成功。
